\chapter{Indirizzamento IP e il protocollo IPv4}

% ══════════════════════════════════════════════════════════════
\section{Il protocollo IP}
% ══════════════════════════════════════════════════════════════

Ricapitolando: il livello di rete deve consegnare i datagrammi \textbf{da host a host}, incapsulando e
decapsulando i segmenti, identificando gli host con un \textbf{indirizzo} e inoltrando i messaggi verso la
destinazione; nella rete i \textbf{router}, con le loro porte, la switching fabric e il processore di controllo che
crea le tabelle, inoltrano i datagrammi ai nodi adiacenti.

\dfn{Internet Protocol (IP)}{
    L'\textbf{Internet Protocol} è il protocollo di livello di rete che assicura la consegna host-a-host. Oggi ne
    esistono due versioni:
    \begin{itemize}
      \item \textbf{IPv4}, la più usata e diffusa;
      \item \textbf{IPv6}, più recente, proposta per sostituire IPv4.
    \end{itemize}
    La funzione più importante di IP è \textbf{identificare gli host} in una rete: l'\textbf{indirizzamento IP} è il
    processo con cui si assegnano gli indirizzi IP ai dispositivi. È cruciale ed è molto complesso su Internet, dove
    ci sono miliardi di host da indirizzare.
}

% ══════════════════════════════════════════════════════════════
\section{Gli indirizzi IPv4}
% ══════════════════════════════════════════════════════════════

Gli indirizzi IPv4 si scrivono di solito in \textbf{notazione decimale puntata} (\emph{dotted-decimal}): ogni byte
dell'indirizzo è scritto in decimale e separato dagli altri da un punto.

\begin{figure}[H]
\centering
\begin{tikzpicture}[every node/.style={font=\small\ttfamily}]
  \foreach \d/\b/\c [count=\i] in {193/11000001/lapp, 32/00100000/ltra, 216/11011000/lnet, 9/00001001/llink} {
    \node[draw, fill=\c, minimum width=2.6cm, minimum height=0.7cm] (b\i) at (\i*2.8,0) {\b};
    \node[font=\large\ttfamily\bfseries] (d\i) at (\i*2.8,1.1) {\d};
    \draw[-{Stealth}] (d\i) -- (b\i);
    \node[font=\scriptsize\sffamily, below=1pt of b\i] {8 bit};
  }
  \foreach \i in {1,2,3} \node[font=\large\bfseries] at (\i*2.8+1.4,1.1) {.};
  \draw[decorate, decoration={brace, amplitude=6pt, mirror}] (1.5,-0.8) -- (12.5,-0.8) node[midway, below=7pt, font=\footnotesize\sffamily] {32 bit = 4 byte};
\end{tikzpicture}
\caption{L'indirizzo 193.32.216.9 in notazione decimale puntata e in binario: ogni numero va da 0 a 255.}
\end{figure}

\begin{itemize}
    \item Ogni dispositivo sull'Internet globale deve avere (in qualche modo) un indirizzo IP \textbf{univoco}.
    \item Poiché un indirizzo è lungo 32 bit, ci sono in tutto $2^{32}$ indirizzi possibili, circa \textbf{4 miliardi}.
\end{itemize}

\subsection{L'indirizzo appartiene all'interfaccia}

Host e dispositivi di rete sono collegati da collegamenti, cablati o senza fili. Un host ha di solito un solo
collegamento verso la rete, mentre un router ne ha molti.

\dfn{Interfaccia}{
    Il confine tra un host (o un router) e un collegamento fisico si chiama \textbf{interfaccia}. Poiché ogni host e
    ogni router può inviare e ricevere datagrammi, IP richiede che \textbf{ogni interfaccia abbia il proprio indirizzo
    IP}. Tecnicamente, quindi, un indirizzo IP è associato all'\textbf{interfaccia}, non all'host o al router che la contiene.
}

Un router con tre interfacce ha tre indirizzi IP; un portatile con scheda cablata e Wi-Fi può averne due.

% ══════════════════════════════════════════════════════════════
\section{Il subnetting}
% ══════════════════════════════════════════════════════════════

Assegnare gli indirizzi alle interfacce non è banale. Assegnarli a caso sarebbe insensato: gli indirizzi non darebbero
alcuna indicazione di dove si trovano gli host, e i router avrebbero bisogno di forwarding table enormi.

L'indirizzamento assomiglia a quello della telefonia tradizionale (prefisso internazionale, prefisso della città,
numero): le reti sono divise \textbf{gerarchicamente} in sottoreti (\textbf{subnet}), con prefissi diversi.

\begin{figure}[H]
\centering
\begin{tikzpicture}[every node/.style={font=\footnotesize\sffamily}]
  \draw[thick, fill=cloudbg] (0,0) circle (2.2cm); \node at (0,1.8) {rete};
  \draw[thick, fill=netblue!20] (-0.9,0.3) circle (0.9cm); \node at (-0.9,0.9) {subnet};
  \draw[thick, fill=netblue!40] (-1.1,0.0) circle (0.35cm); \node[font=\tiny\sffamily] at (-1.1,0.0) {sub-sub};
  \draw[thick, fill=netblue!40] (-0.5,-0.1) circle (0.3cm); \node[font=\tiny\sffamily] at (-0.5,-0.1) {sub-sub};
  \draw[thick, fill=netblue!20] (0.9,0.4) circle (0.7cm); \node at (0.9,0.4) {subnet};
  \draw[thick, fill=netblue!20] (0.2,-1.2) circle (0.7cm); \node at (0.2,-1.2) {subnet};
  \node[text width=7.5cm, anchor=west] at (3,0) {Un indirizzo IP è diviso in due parti:\\[3pt]
    \textbf{parte di rete} (a sinistra): identifica la \textbf{subnet} a cui il nodo è collegato;\\[3pt]
    \textbf{parte di host} (a destra): identifica la singola \textbf{interfaccia} dentro la subnet.\\[3pt]
    Il numero di bit di ciascuna parte \textbf{non è fisso}.};
\end{tikzpicture}
\caption{Reti divise gerarchicamente in subnet.}
\end{figure}

% ══════════════════════════════════════════════════════════════
\section{La subnet mask}
% ══════════════════════════════════════════════════════════════

\dfn{Subnet mask}{
    Per distinguere la parte di subnet dalla parte di host, a un indirizzo IP è associata una \textbf{subnet mask}
    (maschera di sottorete): un numero di 32 bit che ha \textbf{1} nei bit che appartengono alla parte di subnet e
    \textbf{0} in quelli della parte di host.
}

\begin{table}[H]
\centering
\begin{tabular}{lll}
\toprule
 & \textbf{Decimale puntata} & \textbf{Binario} \\
\midrule
Indirizzo IP & 193.32.216.9 & \texttt{\textcolor{netblue}{11000001 00100000 11011000} \textcolor{netred}{00001001}} \\
Subnet mask & 255.255.255.0 & \texttt{\textcolor{netblue}{11111111 11111111 11111111} \textcolor{netred}{00000000}} \\
\midrule
Indirizzo della subnet (AND) & 193.32.216.0 & \texttt{\textcolor{netblue}{11000001 00100000 11011000} 00000000} \\
\bottomrule
\end{tabular}
\caption{Un indirizzo con la sua maschera: in blu la parte di subnet, in rosso la parte di host.}
\end{table}

\dfn{Notazione con lo slash}{
    Un altro modo comune di rappresentare le maschere è la \textbf{notazione con lo slash}: \texttt{193.32.216.0/24}
    indica l'indirizzo della subnet, dove \texttt{/24} dice che i \textbf{24 bit più a sinistra} dell'indirizzo sono
    dedicati alla subnet (maschera 255.255.255.0).
}

\begin{esempio}{Tutti gli indirizzi di una subnet}
La subnet \texttt{193.32.216.0/24} contiene tutti gli indirizzi che hanno i primi 24 bit uguali:
\begin{itemize}
    \item il più piccolo: \texttt{11000001 00100000 11011000 00000000} = \texttt{193.32.216.0};
    \item il più grande: \texttt{11000001 00100000 11011000 11111111} = \texttt{193.32.216.255}.
\end{itemize}
Sono $2^{32-24} = 2^8 = 256$ indirizzi. Due sono riservati: quello con la parte di host tutta a 0 indica la
\textbf{rete stessa} (193.32.216.0), quello con la parte di host tutta a 1 è l'indirizzo di \textbf{broadcast}
(193.32.216.255). Restano \textbf{254} indirizzi assegnabili alle interfacce.
\end{esempio}

\warning{Quanti host in una subnet /x}{
    Una subnet \texttt{/x} contiene $2^{32-x}$ indirizzi, di cui $2^{32-x} - 2$ assegnabili agli host.
    Esempi: /24 $\rightarrow$ 254 host; /25 $\rightarrow$ 126; /26 $\rightarrow$ 62; /30 $\rightarrow$ 2 (usata per i
    collegamenti punto-punto tra router); /16 $\rightarrow$ 65\,534.
}

\subsection{Quali maschere sono valide}

Una subnet mask identifica un \textbf{prefisso}: gli 1 devono quindi essere tutti \textbf{a sinistra} e
\textbf{contigui}.

\begin{table}[H]
\centering
\begin{tabular}{llc}
\toprule
\textbf{Decimale puntata} & \textbf{Binario} & \textbf{Valida?} \\
\midrule
255.255.255.0 & \texttt{11111111 11111111 11111111 00000000} & \itick\ /24 \\
255.255.128.0 & \texttt{11111111 11111111 10000000 00000000} & \itick\ /17 \\
255.224.0.0 & \texttt{11111111 11100000 00000000 00000000} & \itick\ /11 \\
255.255.10.0 & \texttt{11111111 11111111 00001010 00000000} & \icross\ 1 non contigui \\
63.255.255.0 & \texttt{00111111 11111111 11111111 00000000} & \icross\ non parte da sinistra \\
\bottomrule
\end{tabular}
\caption{Maschere valide e non valide.}
\end{table}

\info{I valori possibili di un byte della maschera}{
    Poiché gli 1 sono contigui, ogni byte di una maschera può valere solo: 0, 128, 192, 224, 240, 248, 252, 254, 255.
}

\subsection{Verificare se due host sono nella stessa subnet}

Per capire se due host sono nella stessa subnet basta controllare se i loro \textbf{prefissi} sono uguali, cioè se
\textbf{IP1 AND maschera} = \textbf{IP2 AND maschera}.

\begin{table}[H]
\centering
\begin{tabular}{lllc}
\toprule
\textbf{IP 1} & \textbf{Maschera} & \textbf{IP 2} & \textbf{Esito} \\
\midrule
231.23.11.117 & 255.255.255.0 & 231.23.11.9 & stessa subnet \\
10.54.32.1 & 255.255.128.0 & 10.54.60.203 & stessa subnet \\
110.32.100.10 & 255.224.0.0 & 110.64.100.11 & subnet diverse \\
121.50.79.4 & 255.0.0.0 & 121.10.36.240 & stessa subnet \\
121.50.79.4 & 255.0.0.0 & 120.50.79.5 & subnet diverse \\
\bottomrule
\end{tabular}
\caption{Esempi di confronto tra indirizzi.}
\end{table}

\begin{esempio}{Due casi svolti in binario}
\textbf{10.54.32.1 e 10.54.60.203 con maschera /17.} Il terzo byte della maschera è 128 = \texttt{10000000}: conta solo il suo
primo bit.
\begin{itemize}
    \item $32 = \texttt{\textcolor{netblue}{0}0100000}$, $60 = \texttt{\textcolor{netblue}{0}0111100}$: il primo bit è 0 in entrambi, e i primi due byte coincidono (10.54) $\Rightarrow$ \textbf{stessa subnet} (10.54.0.0/17).
\end{itemize}
\textbf{110.32.100.10 e 110.64.100.11 con maschera /11.} Il secondo byte della maschera è 224 = \texttt{11100000}: contano i primi 3 bit.
\begin{itemize}
    \item $32 = \texttt{\textcolor{netblue}{001}00000}$, $64 = \texttt{\textcolor{netred}{010}00000}$: i primi tre bit sono diversi $\Rightarrow$ \textbf{subnet diverse}.
\end{itemize}
\end{esempio}

Host di subnet diverse possono comunque comunicare, ma gli amministratori di rete possono decidere di trattare
diversamente i messaggi che entrano o escono da una subnet (per esempio bloccandoli o filtrandoli).

% ══════════════════════════════════════════════════════════════
\section{Router e subnet}
% ══════════════════════════════════════════════════════════════

Il dispositivo che collega subnet diverse è il \textbf{router}. Per definizione un router ha più porte fisiche,
che possono corrispondere a interfacce di rete diverse con indirizzi IP diversi. I router stanno spesso sul
\textbf{bordo} di più subnet e gestiscono i messaggi in entrata e in uscita: si dice che fanno da
\textbf{gateway}. Grazie a questa posizione privilegiata offrono spesso altri servizi oltre al routing:

\begin{itemize}
    \item \textbf{IP masquerading} (NAT);
    \item \textbf{assegnazione dinamica degli indirizzi} (DHCP);
    \item \textbf{protezione} (firewall).
\end{itemize}

\begin{esempio}{Tre subnet e un router}
\begin{center}
\includegraphics[width=0.5\linewidth]{cap14/3subnet.png}
\end{center}
Un router con tre interfacce collega sette host, divisi in tre subnet (sinistra, destra, sotto), ciascuna collegata a
un'interfaccia del router.
\begin{itemize}
    \item La subnet di sinistra ha indirizzo \texttt{223.1.1.0/24}: tutte le interfacce hanno indirizzi della forma
          \texttt{223.1.1.xxx} --- \texttt{223.1.1.1}, \texttt{223.1.1.2}, \texttt{223.1.1.3} (host) e \texttt{223.1.1.4} (router).
    \item Le altre due subnet sono \texttt{223.1.2.0/24} e \texttt{223.1.3.0/24}.
\end{itemize}
\end{esempio}

\begin{esempio}{Tre router e sei subnet}
\begin{center}
\includegraphics[width=0.42\linewidth]{cap14/6subnet.png}
\end{center}
Tre router sono collegati tra loro da collegamenti diretti punto-punto. Ogni router ha tre interfacce: una per ciascun
collegamento punto-punto e una per il collegamento broadcast verso una coppia di host. Le subnet sono \textbf{sei}:
\begin{itemize}
    \item \texttt{223.1.1.0/24} (R1--host), \texttt{223.1.2.0/24} (R2--host), \texttt{223.1.3.0/24} (R3--host);
    \item \texttt{223.1.9.0/24} (R1--R2), \texttt{223.1.8.0/24} (R2--R3), \texttt{223.1.7.0/24} (R3--R1).
\end{itemize}
Per contare le subnet si ``staccano'' le interfacce da host e router: ogni isola di rete che resta è una subnet, e anche un
collegamento punto-punto tra due router è una subnet. Le organizzazioni medio-grandi (aziende, università) hanno tipicamente
molte subnet interconnesse.
\end{esempio}

\begin{figure}[H]
\centering
\begin{tikzpicture}[every node/.style={font=\footnotesize\sffamily}]
  \node (pc) at (0,0) {\icon[1cm]{laptop}};
  \node[below=-2pt of pc, align=center] {host\\\texttt{192.168.1.4}\\maschera \texttt{255.255.255.0}};
  \node (r) at (5,0) {\icon[1.3cm]{accesspoint}};
  \node[below=-2pt of r, align=center] {router di casa (gateway)\\\texttt{192.168.1.1}};
  \node[netcloud, minimum width=2.4cm] (i) at (9.5,0) {Internet};
  \draw[wireless] (pc) -- (r); \draw[wired, line width=2pt] (r) -- (i);
  \node[draw, rounded corners, fill=cloudbg, text width=6cm, align=left, anchor=north] at (5,-1.8) {Per un indirizzo \textbf{nella stessa subnet} (192.168.1.x) l'host consegna direttamente; per \textbf{tutti gli altri} invia i pacchetti al \textbf{gateway}, cioè al router.};
\end{tikzpicture}
\caption{La configurazione tipica di un host di casa: indirizzo, maschera e gateway predefinito.}
\end{figure}

\subsection{ifconfig e ip}

Su Linux si può controllare il proprio indirizzo con \texttt{ifconfig} (\emph{interface configuration}; l'equivalente su
Windows è \texttt{ipconfig}), che elenca tutte le interfacce di rete della macchina con la loro configurazione
(indirizzi, maschere, \dots). Il comando moderno è \texttt{ip}.

\begin{lstlisting}[language=bash]
$ ifconfig          # configurazione di tutte le interfacce
$ ip addr           # lo stesso, con il comando moderno
\end{lstlisting}

% ══════════════════════════════════════════════════════════════
\section{L'indirizzamento su Internet}
% ══════════════════════════════════════════════════════════════

Dividere le grandi reti in reti più piccole è particolarmente importante su Internet, dove vanno collegati miliardi di
dispositivi (e ancora più interfacce). Gli indirizzi vanno assegnati con cura per evitare che:

\begin{itemize}
    \item le forwarding table dei router diventino enormi;
    \item interfacce diverse abbiano lo stesso indirizzo;
    \item si esauriscano gli indirizzi.
\end{itemize}

L'approccio è dare alle organizzazioni (ISP, aziende, enti) delle subnet, in due modi: l'indirizzamento
\textbf{classful} (vecchio, non più usato in pratica) e quello \textbf{classless} (attuale).

\subsection{Indirizzamento classful}

\begin{table}[H]
\centering
\begin{tabular}{llllr}
\toprule
\textbf{Classe} & \textbf{Formato} & \textbf{Primo byte} & \textbf{Esempio} & \textbf{Indirizzi per rete} \\
\midrule
A & a.b.c.d/8 & 0--127 & 10.X.X.X & oltre 16 milioni \\
B & a.b.c.d/16 & 128--191 & 150.10.X.X & 65\,536 \\
C & a.b.c.d/24 & 192--223 & 200.10.10.X & 256 (254 host) \\
\bottomrule
\end{tabular}
\caption{Le classi di indirizzi (le classi D, 224--239, ed E, 240--255, erano riservate a multicast e usi sperimentali).}
\end{table}

\begin{esempio}{Lo spreco del classful}
Un'organizzazione che ha bisogno di 300 indirizzi non può usare una classe C (troppo piccola): riceve una classe B, per
esempio \texttt{241.115.0.0/16}, con tutti i suoi 65\,536 indirizzi. Ne usa 300, e ne \textbf{spreca oltre 65\,000}.
\end{esempio}

\subsection{Indirizzamento classless (CIDR)}

\dfn{CIDR}{
    L'approccio moderno, più flessibile, si chiama \textbf{Classless InterDomain Routing} (CIDR, pronunciato come
    ``\emph{cider}''). A un'organizzazione si può assegnare un indirizzo di rete della forma \texttt{a.b.c.d/x}, con
    $x$ \textbf{qualsiasi}.
    \begin{itemize}
      \item La parte di rete dell'indirizzo si chiama \textbf{prefisso}.
      \item L'insieme di indirizzi riservati all'organizzazione si chiama \textbf{blocco}.
    \end{itemize}
}

\begin{esempio}{300 indirizzi con CIDR}
Servono 300 indirizzi: bastano 9 bit di host ($2^9 = 512 \geq 300$), quindi un prefisso di $32 - 9 = 23$ bit. L'organizzazione
riceve per esempio \texttt{241.115.2.0/23}: 512 indirizzi, con soli 212 sprecati (invece di oltre 65\,000).
\end{esempio}

\warning{Blocchi che si sovrappongono}{
    I blocchi possono sovrapporsi: in questo caso è la \textbf{lunghezza del prefisso} a distinguerli. Per esempio
    \texttt{10.10.10.15/16} \textbf{non} è nello stesso blocco di \texttt{10.10.10.14/24}: il primo appartiene a
    \texttt{10.10.0.0/16}, il secondo a \texttt{10.10.10.0/24}.
}

\subsection{Subnet interne a un blocco}

Dentro un blocco CIDR si possono creare \textbf{subnet interne} all'organizzazione. Il blocco \texttt{241.115.0.0/16}, per
esempio, si può dividere in:
\texttt{241.115.1.0/24} (subnet 1), \texttt{241.115.2.0/24} (subnet 2), \texttt{241.115.3.0/24} (subnet 3), e così via.

\begin{figure}[H]
\centering
\begin{tikzpicture}[every node/.style={font=\small\ttfamily}]
  \node[draw, fill=lnet!60, minimum width=4.4cm, minimum height=0.7cm] (p) at (0,0) {11110001 01110011};
  \node[draw, fill=ltra!60, minimum width=2.2cm, minimum height=0.7cm, right=0pt of p] (s) {SSSSSSSS};
  \node[draw, fill=lapp!60, minimum width=2.2cm, minimum height=0.7cm, right=0pt of s] (h) {HHHHHHHH};
  \draw[decorate, decoration={brace, amplitude=5pt}] (p.north west) -- (p.north east) node[midway, above=6pt, font=\footnotesize\sffamily] {prefisso dell'organizzazione (/16)};
  \draw[decorate, decoration={brace, amplitude=5pt}] (s.north west) -- (s.north east) node[midway, above=6pt, font=\footnotesize\sffamily] {subnet};
  \draw[decorate, decoration={brace, amplitude=5pt}] (h.north west) -- (h.north east) node[midway, above=6pt, font=\footnotesize\sffamily] {host};
\end{tikzpicture}
\caption{Un blocco /16 diviso internamente in subnet /24.}
\end{figure}

\begin{esempio}{Dividere un blocco in sottoblocchi uguali}
Un ISP ha il blocco \texttt{200.23.16.0/20} e vuole dividerlo tra 8 organizzazioni. Servono $\log_2 8 = 3$ bit in più: i
sottoblocchi sono /23, ciascuno da $2^9 = 512$ indirizzi. Il terzo byte vale $16 = \texttt{0001}\,\texttt{0000}$: i 4 bit
evidenziati sono del prefisso /20, i 3 successivi numerano l'organizzazione.
\begin{center}
\begin{tabular}{lll}
\toprule
\textbf{Organizzazione} & \textbf{Terzo byte} & \textbf{Blocco} \\
\midrule
0 & \texttt{0001 000}0 & 200.23.16.0/23 \\
1 & \texttt{0001 001}0 & 200.23.18.0/23 \\
2 & \texttt{0001 010}0 & 200.23.20.0/23 \\
\dots & \dots & \dots \\
7 & \texttt{0001 111}0 & 200.23.30.0/23 \\
\bottomrule
\end{tabular}
\end{center}
\end{esempio}

% ══════════════════════════════════════════════════════════════
\section{L'aggregazione degli indirizzi}
% ══════════════════════════════════════════════════════════════

L'indirizzamento basato sui prefissi è utilissimo ai dispositivi che collegano prefissi diversi: un router può
ricordare (cioè scrivere nella forwarding table) solo i \textbf{prefissi}. Quando riceve un messaggio con un certo
prefisso lo inoltra verso un router più specifico, e così via.

\dfn{Aggregazione degli indirizzi}{
    L'\textbf{aggregazione degli indirizzi} (\emph{address aggregation} o \emph{route summarization}) consiste
    nell'annunciare al resto della rete un \textbf{unico prefisso} che riassume molti blocchi più piccoli.
}

\begin{figure}[H]
\centering
\begin{tikzpicture}[every node/.style={font=\scriptsize\sffamily}, org/.style={draw, rounded corners, fill=cloudbg, minimum width=3.4cm, align=center}]
  \node[org] (o0) at (0,2.4) {Organizzazione 0\\\texttt{200.23.16.0/23}};
  \node[org] (o1) at (0,1.2) {Organizzazione 1\\\texttt{200.23.18.0/23}};
  \node[org] (o2) at (0,0) {Organizzazione 2\\\texttt{200.23.20.0/23}};
  \node at (0,-0.8) {$\vdots$};
  \node[org] (o7) at (0,-1.7) {Organizzazione 7\\\texttt{200.23.30.0/23}};
  \node (isp) at (5,0.5) {\icon[1.1cm]{router}}; \node[below=-2pt of isp] {ISP1};
  \node[netcloud, minimum width=2.4cm] (net) at (9.6,0.5) {Internet};
  \foreach \o in {o0,o1,o2,o7} \draw[thick] (\o.east) -- (isp);
  \draw[wired, line width=2pt] (isp) -- (net);
  \node[draw, fill=llink!60, rounded corners, text width=4cm, align=center] at (7.3,2.2) {``mandatemi tutto ciò che inizia con \texttt{200.23.16.0/20}''};
\end{tikzpicture}
\caption{ISP1 annuncia un solo prefisso /20 per le otto organizzazioni: il resto di Internet ha bisogno di una sola riga.}
\end{figure}

Che succede se l'Organizzazione 1 cambia fornitore e passa a ISP2, portandosi dietro il proprio blocco?

\begin{figure}[H]
\centering
\begin{tikzpicture}[every node/.style={font=\scriptsize\sffamily}, org/.style={draw, rounded corners, fill=cloudbg, minimum width=3.4cm, align=center}]
  \node[org] (o0) at (0,2.4) {Organizzazione 0\\\texttt{200.23.16.0/23}};
  \node[org] (o7) at (0,1.0) {Organizzazione 7\\\texttt{200.23.30.0/23}};
  \node[org, fill=netorange!20] (o1) at (0,-1.2) {Organizzazione 1\\\texttt{200.23.18.0/23}};
  \node (isp1) at (5,1.7) {\icon[1cm]{router}}; \node[below=-2pt of isp1] {ISP1};
  \node (isp2) at (5,-1.2) {\icon[1cm]{router}}; \node[below=-2pt of isp2] {ISP2};
  \node[netcloud, minimum width=2.4cm] (net) at (9.6,0.3) {Internet};
  \draw[thick] (o0.east) -- (isp1); \draw[thick] (o7.east) -- (isp1); \draw[thick, netorange] (o1.east) -- (isp2);
  \draw[wired, line width=2pt] (isp1) -- (net); \draw[wired, line width=2pt] (isp2) -- (net);
  \node[draw, fill=llink!60, rounded corners, text width=3.6cm, align=center] at (7.6,2.6) {``tutto ciò che inizia con \texttt{200.23.16.0/20}''};
  \node[draw, fill=netorange!20, rounded corners, text width=3.6cm, align=center] at (7.6,-2.0) {``\dots ma \texttt{200.23.18.0/23} mandatelo a me!''};
\end{tikzpicture}
\caption{ISP2 annuncia il prefisso più specifico: grazie alla longest prefix matching i pacchetti per l'Organizzazione 1 vanno a ISP2, tutti gli altri del /20 a ISP1.}
\end{figure}

\begin{itemize}
    \item Un indirizzo dell'Organizzazione 1 corrisponde sia al /20 di ISP1 sia al /23 di ISP2; per la \textbf{longest prefix
          matching} vince il /23: nessuna ambiguità.
    \item Blocchi e sottoblocchi possono quindi essere assegnati a linee diverse senza ambiguità, con grande flessibilità.
    \item L'aggregazione è tanto più efficace quanto più i blocchi restano \textbf{raggruppati} (clusterizzati): ogni
          eccezione è una riga in più nelle tabelle.
\end{itemize}

% ══════════════════════════════════════════════════════════════
\section{Come ottenere un blocco di indirizzi}
% ══════════════════════════════════════════════════════════════

Per ottenere un blocco di indirizzi per le subnet di un'organizzazione ci sono due strade:

\begin{enumerate}
    \item \textbf{Contattare un ISP}, che fornisce indirizzi presi da un blocco più grande che gli è già stato assegnato. Per
          esempio l'ISP può avere il blocco \texttt{200.23.16.0/20} e dividerlo in sottoblocchi di dimensione variabile
          secondo le proprie politiche.
    \item \textbf{Rivolgersi all'ICANN} (\emph{Internet Corporation for Assigned Names and Numbers}), l'autorità globale no-profit
          responsabile dello spazio degli indirizzi IP (assegna i blocchi agli ISP, tramite registri regionali). L'ICANN gestisce
          anche i root server DNS, assegna i nomi di dominio e risolve le controversie sui nomi.
\end{enumerate}

% ══════════════════════════════════════════════════════════════
\section{Il datagramma IPv4}
% ══════════════════════════════════════════════════════════════

Il pacchetto del livello di rete di Internet si chiama \textbf{datagramma} (come in UDP).

\begin{figure}[H]
\centering
\begin{tikzpicture}[every node/.style={font=\scriptsize\sffamily}, x=0.45cm, y=0.62cm]
  \node[field, fill=lnet!60, minimum width=1.8cm, minimum height=0.62cm, anchor=north west] at (0,0) {versione};
  \node[field, fill=lnet!60, minimum width=1.8cm, minimum height=0.62cm, anchor=north west] at (4,0) {lungh. header};
  \node[field, fill=lnet!40, minimum width=3.6cm, minimum height=0.62cm, anchor=north west] at (8,0) {tipo di servizio};
  \node[field, fill=lnet!40, minimum width=7.2cm, minimum height=0.62cm, anchor=north west] at (16,0) {lunghezza del datagramma (byte)};
  \node[field, fill=netorange!30, minimum width=7.2cm, minimum height=0.62cm, anchor=north west] at (0,-1) {identificatore (16)};
  \node[field, fill=netorange!30, minimum width=1.35cm, minimum height=0.62cm, anchor=north west] at (16,-1) {flag (3)};
  \node[field, fill=netorange!30, minimum width=5.85cm, minimum height=0.62cm, anchor=north west] at (19,-1) {offset del frammento (13)};
  \node[field, fill=netred!20, minimum width=3.6cm, minimum height=0.62cm, anchor=north west] at (0,-2) {time to live (8)};
  \node[field, fill=ltra!50, minimum width=3.6cm, minimum height=0.62cm, anchor=north west] at (8,-2) {protocollo superiore};
  \node[field, fill=lnet!40, minimum width=7.2cm, minimum height=0.62cm, anchor=north west] at (16,-2) {checksum dell'header (16)};
  \node[field, fill=llink!70, minimum width=14.4cm, minimum height=0.62cm, anchor=north west] at (0,-3) {indirizzo IP sorgente (32)};
  \node[field, fill=llink!70, minimum width=14.4cm, minimum height=0.62cm, anchor=north west] at (0,-4) {indirizzo IP destinazione (32)};
  \node[field, fill=lphy!60, minimum width=14.4cm, minimum height=0.62cm, anchor=north west] at (0,-5) {opzioni (se presenti)};
  \node[field, fill=cloudbg, minimum width=14.4cm, minimum height=1.2cm, anchor=north west] at (0,-6) {dati (tipicamente un segmento TCP o UDP)};
  \draw[{Stealth}-{Stealth}] (0,0.5) -- node[above] {32 bit} (32,0.5);
  \draw[decorate, decoration={brace, amplitude=5pt}] (32.4,0) -- (32.4,-5) node[midway, right=6pt, align=left] {header:\\20 byte\\senza opzioni};
\end{tikzpicture}
\caption{Il datagramma IPv4.}
\end{figure}

\begin{itemize}
    \item \textbf{Versione} (4 bit): la versione di IP (4 o 6); il formato dipende dalla versione.
    \item \textbf{Lunghezza dell'header} (4 bit): la dimensione dell'header, non fissa per via delle opzioni; serve a sapere
          dove iniziano i dati. Senza opzioni l'header è di \textbf{20 byte}.
    \item \textbf{Tipo di servizio} (TOS, 8 bit): proprietà specifiche del datagramma (per esempio tempo reale o no),
          definite dagli amministratori dei router.
    \item \textbf{Lunghezza del datagramma} (16 bit): lunghezza in byte di header + dati. I datagrammi raramente superano
          i 1500 byte, sui 65\,535 possibili.
    \item \textbf{Identificatore} (16 bit): numero progressivo che identifica il datagramma (usato nella frammentazione).
    \item \textbf{Flag e offset del frammento} (3 + 13 bit): proprietà e posizione del frammento (usati nella frammentazione).
    \item \textbf{Time-to-live} (TTL, 8 bit): impedisce che i datagrammi circolino per sempre. Ogni router che elabora il
          datagramma lo \textbf{decrementa di 1}; se arriva a 0, il router \textbf{scarta} il datagramma.
    \item \textbf{Protocollo superiore} (8 bit): il protocollo di trasporto a cui consegnare i dati (per esempio \textbf{6} per
          TCP, \textbf{17} per UDP; l'elenco completo è sul sito della IANA).
    \item \textbf{Checksum dell'header} (16 bit): complemento a 1 della somma a 16 bit dell'header. I router lo controllano e di
          solito scartano i datagrammi errati. Va \textbf{ricalcolato a ogni router}, perché TTL e opzioni possono cambiare.
    \item \textbf{Indirizzi IP sorgente e destinazione} (32 + 32 bit).
    \item \textbf{Opzioni} (lunghezza variabile): estendono l'header con funzionalità aggiuntive. Aggiungono complessità (la
          dimensione non è nota a priori) e in IPv6 non esistono più.
    \item \textbf{Dati} (lunghezza variabile): il payload, tipicamente un segmento TCP o UDP.
\end{itemize}

\info{Quanto pesano gli header}{
    Un segmento TCP dentro un datagramma IP porta almeno 40 byte di header (20 TCP + 20 IP), più l'header del
    livello applicativo. Per messaggi piccoli l'overhead è significativo.
}

% ══════════════════════════════════════════════════════════════
\section{La frammentazione IPv4}
% ══════════════════════════════════════════════════════════════

Il primo problema dei datagrammi IP è che possono dover essere \textbf{frammentati} a causa del livello di
collegamento: protocolli di collegamento diversi trasportano pacchetti di dimensioni diverse (i frame Ethernet, per
esempio, portano al massimo 1500 byte di dati).

\dfn{MTU}{
    I datagrammi IP vengono incapsulati in frame di collegamento per essere trasportati da un nodo all'altro. La
    quantità massima di dati che un frame può trasportare è la \textbf{Maximum Transmission Unit} (MTU). A livello di
    collegamento la lunghezza del frame è limitata dai vincoli della trasmissione fisica.
}

\begin{figure}[H]
    \centering
    \includegraphics[width=0.55\linewidth]{cap14/frammentazione.png}
    \caption{Un router frammenta un datagramma da 4000 byte per un collegamento con MTU di 1500 byte; la destinazione riassembla i tre frammenti.}
\end{figure}

Un router che collega due collegamenti con MTU diverse può ricevere dal collegamento in ingresso un datagramma che non
entra in quello in uscita. Deve allora dividere il payload in due o più datagrammi più piccoli, i
\textbf{frammenti}, che vengono incapsulati in frame e inoltrati.

\subsection{Come funziona la frammentazione}

Quando un router frammenta un datagramma:

\begin{itemize}
    \item copia nei nuovi frammenti lo stesso \textbf{identificatore}, lo stesso indirizzo sorgente e lo stesso indirizzo di destinazione;
    \item imposta il campo \textbf{offset} di ciascun frammento con numeri progressivi (la posizione dei suoi dati nel
          datagramma originale, misurata in multipli di 8 byte);
    \item imposta il \textbf{flag} (``more fragments'') a 1 in tutti i frammenti tranne l'ultimo, dove vale 0: così si capisce
          che i frammenti sono finiti.
\end{itemize}

\begin{esempio}{4000 byte su un collegamento con MTU di 1500}
Un datagramma di 4000 byte (20 di header + 3980 di dati) con identificatore 777 deve attraversare un collegamento con MTU
di 1500 byte. Ogni frammento può portare al massimo $1500 - 20 = 1480$ byte di dati (multiplo di 8).
\begin{center}
\begin{tabular}{ccccc}
\toprule
\textbf{Frammento} & \textbf{Byte di dati} & \textbf{Identificatore} & \textbf{Offset} & \textbf{Flag} \\
\midrule
1 & 1480 (byte 0--1479) & 777 & $0/8 = 0$ & 1 \\
2 & 1480 (byte 1480--2959) & 777 & $1480/8 = 185$ & 1 \\
3 & 1020 (byte 2960--3979) & 777 & $2960/8 = 370$ & 0 \\
\bottomrule
\end{tabular}
\end{center}
Il terzo frammento è lungo $1020 + 20 = 1040$ byte in tutto. La somma dei dati è $1480 + 1480 + 1020 = 3980$ byte.
\end{esempio}

\warning{Il riassemblaggio avviene negli host}{
    I frammenti vanno riassemblati prima di arrivare al trasporto, perché TCP e UDP si aspettano segmenti completi. I
    progettisti di IPv4 ritennero che riassemblare nei router avrebbe complicato il protocollo e rallentato i router:
    in IPv4 il \textbf{riassemblaggio è eseguito dagli end system}.
    \begin{itemize}
      \item Se l'host riceve più datagrammi con gli stessi indirizzi e lo stesso identificatore, capisce che il datagramma
            originale è stato frammentato.
      \item Deve ricostruire il datagramma originale dai frammenti, usando offset e flag. Se un frammento si perde, si
            scarta l'intero datagramma.
    \end{itemize}
}

% ══════════════════════════════════════════════════════════════
\section{Riepilogo}
% ══════════════════════════════════════════════════════════════

\riepilogo{
\begin{itemize}
    \item Un indirizzo IPv4 è di 32 bit, scritto in notazione decimale puntata; appartiene a un'\textbf{interfaccia}.
    \item Indirizzo = \textbf{prefisso} (subnet) + parte di host; la \textbf{subnet mask} (o /x) dice quanti bit sono di prefisso.
          Due host sono nella stessa subnet se IP AND maschera coincide.
    \item Una subnet /x ha $2^{32-x}$ indirizzi e $2^{32-x}-2$ host.
    \item Dal \textbf{classful} (A, B, C, sprechi) al \textbf{CIDR} (a.b.c.d/x qualsiasi); \textbf{aggregazione} dei prefissi con
          longest prefix matching per le eccezioni.
    \item Header IPv4: 20 byte + opzioni; campi chiave TTL, protocollo, checksum, indirizzi, identificatore/flag/offset.
    \item \textbf{Frammentazione}: se il datagramma supera l'MTU, i router lo dividono; gli host lo riassemblano.
\end{itemize}
}
